\section{Discussion}
\label{sec:discussion}

The central V1 result is negative for universal deep-fusion superiority. A
simple, controlled classical feature concatenation method (C4) produced the
lowest clean ROCCH-EER and retained first rank under every predeclared image
degradation. The compact deep feature model D2 was the strongest neural family
and exhibited a favorable robustness-loss point estimate, yet did not close the
clean discrimination gap to C4. Likewise, D1 did not beat the strongest
classical score fusion, while D3S only narrowly improved on its quality-aware
classical counterpart. These outcomes are retained because they directly answer
Q1--Q3 under matched final trials and frozen selection decisions.

The Q2 frontier explains why a single scalar ranking would be misleading.
Feature methods bought lower error at materially higher fusion cost, whereas
small score combiners remained non-dominated through latency and parameter
economy. Frontier membership is descriptive for the four predeclared dimensions
and is not evidence that a method would dominate for a particular operational
utility or target device.

The Q3 findings separate winner stability from full-ranking stability. C4's
first position survived every locked blur and contrast condition, but the
ordering among the remaining families was condition-dependent. Under complete
loss of one modality, the frozen fallback intentionally eliminates method
differences. This is operationally predictable behavior, but it cannot support
a claim that the learned gates exploit missingness better than classical
systems.

\subsection{Limitations and V2 boundary}

V1 uses 800 images associated with 20 public biometric-instance identifiers,
not the complete 33,600-image NUPT-FPV archive. Public documentation does not
establish how these identifiers map to independent human volunteers. The dense
all-versus-all impostor construction also makes a one-way identity bootstrap an
inadequate substitute for person-level inference. For these reasons V1 reports
point estimates only; the results cannot be generalized to a human population,
the full database, other sensors, demographic groups, presentation attacks, or
deployment conditions.

Only the frozen blur, contrast, and complete-modality-absence conditions were
tested. Noise, compression, occlusion, cross-dataset transfer, security,
fairness, and end-to-end encoder variation remain outside the confirmatory V1
scope. Fusion-only timings were measured under the recorded workflow
environment and should not be treated as device-independent latency claims.

V2 should preserve the Q1--Q3 comparison and acquire the complete archive with
verified person-to-finger mapping. Person-disjoint partitions and an appropriate
dependence-aware inferential design can then quantify uncertainty and test
whether the observed classical advantage and condition-wise winner stability
replicate at population-relevant scale. V2 is a replication/generalization
stage, not an opportunity to retune V1 after observing its final scores.
